
\subsection{Anwendungsbeispiel Palettieren-LEA}
\label{subsec:PalettierenLea}

Zur Veranschaulichung der in Kapitel~\ref{chap:Art1LEA} beschriebenen Konzepte zeigt dieser Abschnitt die beispielhafte Anwendung dieser Konzepte auf ein Palettieren-LEA (PAL), welches einer CES-basierten Arbeitsweise folgt.

\subsubsection{Dienstespezifikation}

Tabelle~\ref{tab:PalDienst} gibt einen Überblick über den Dienst des PAL inkl. dessen Prozeduren, Parametern, Reportwerten und Prozesswerten.

% Steckbrief PAL-Dienst
\begin{footnotesize}
	\begin{longtable}{|
		>{\columncolor[HTML]{E0E0E0}}l
		>{\columncolor[HTML]{E0E0E0}}l
		>{\columncolor[HTML]{E0E0E0}}l 
		>{\columncolor[HTML]{FFFFFF}}l |}
		\caption{Steckbrief eines beispielhaften Palettierendienstes}
		\label{tab:PalDienst}\\
		\hline

		\multicolumn{4}{|L{14.8cm}|}{\cellcolor[HTML]{E0E0E0}\textbf{Palettieren-Dienst}} 
		\\ \hline

		\multicolumn{4}{|L{14.8cm}|}{\cellcolor[HTML]{FFFFFF} Dieser Dienst dient zur Palettierung von Säcken auf Paletten nach einem festgelegten Packmuster. Die Säcke werden dabei gemäß bestimmter Auftragsdaten immer gleichartig verarbeitet. } 
		\\ \hline

		\multicolumn{4}{|L{14.8cm}|}{\cellcolor[HTML]{E0E0E0}\textbf{Prozeduren}} 
		\\ \hline

		\multicolumn{1}{|L{2.5cm}|}{\cellcolor[HTML]{E0E0E0}\textbf{Prozedur-ID}}        	& 
		\multicolumn{1}{L{3.3cm}|}{\cellcolor[HTML]{E0E0E0}\textbf{Name}} 			& 
		\multicolumn{2}{L{8.65cm}|}{\cellcolor[HTML]{E0E0E0}\textbf{Beschreibung}} 
		\\ \hline

		\multicolumn{1}{|L{2.5cm}|}{\cellcolor[HTML]{FFFFFF} 16\#1}                     & 
		\multicolumn{1}{L{3.3cm}|}{\cellcolor[HTML]{FFFFFF} CES\_Continuous (kontinuierlich)}              &                                              
		\multicolumn{2}{L{8.65cm}|}{\cellcolor[HTML]{FFFFFF} Kontinuierliches Palettieren von Säcken mit immer gleichem Packmuster}              
		\\ \hline

		\multicolumn{1}{|L{2.5cm}|}{\cellcolor[HTML]{FFFFFF} 16\#2}                     & 
		\multicolumn{1}{L{3.3cm}|}{\cellcolor[HTML]{FFFFFF} CES\_BagQuantity (selbstbeendend)}              &                                              
		\multicolumn{2}{L{8.65cm}|}{\cellcolor[HTML]{FFFFFF} Palettieren einer bestimmten Anzahl an Säcken mit immer gleichem Packmuster}              
		\\ \hline

		\multicolumn{1}{|L{2.5cm}|}{\cellcolor[HTML]{FFFFFF} 16\#3}                     & 
		\multicolumn{1}{L{3.3cm}|}{\cellcolor[HTML]{FFFFFF} CES\_PalletQuantity (selbstbeendend)}              &                                              
		\multicolumn{2}{L{8.65cm}|}{\cellcolor[HTML]{FFFFFF} Palettieren einer bestimmten Anzahl an Paletten mit immer gleichem Packmuster}              
		\\ \hline

		\multicolumn{4}{|L{14.8cm}|}{\cellcolor[HTML]{E0E0E0}\textbf{Auftragsspezifische Parameter (Prozedurparameter)}} 
		\\ \hline

		\multicolumn{1}{|L{2.5cm}|}{\cellcolor[HTML]{E0E0E0}\textbf{Name}}        	& 
		\multicolumn{1}{L{3.3cm}|}{\cellcolor[HTML]{E0E0E0}\textbf{Schnittstellentyp}} 			& 
		\multicolumn{1}{L{3.17cm}|}{\cellcolor[HTML]{E0E0E0}\textbf{Prozedur}} 			& 
		\multicolumn{1}{L{5.3cm}|}{\cellcolor[HTML]{E0E0E0}\textbf{Beschreibung}} 
		\\ \hline

		
		\multicolumn{1}{|L{2.5cm}|}{\cellcolor[HTML]{FFFFFF} ProductionId}        	& 
		\multicolumn{1}{L{3.3cm}|}{\cellcolor[HTML]{FFFFFF} StringServParam} 			& 
		\multicolumn{1}{L{3.17cm}|}{\cellcolor[HTML]{FFFFFF} 16\#1, 16\#2, 16\#3} 			& 
		\multicolumn{1}{L{5.3cm}|}{\cellcolor[HTML]{FFFFFF} Auftragsnummer} 
		\\ \hline

		\multicolumn{1}{|L{2.5cm}|}{\cellcolor[HTML]{FFFFFF} ProductId}        	& 
		\multicolumn{1}{L{3.3cm}|}{\cellcolor[HTML]{FFFFFF} DIntServParam} 			& 
		\multicolumn{1}{L{3.17cm}|}{\cellcolor[HTML]{FFFFFF} 16\#1, 16\#2, 16\#3} 			& 
		\multicolumn{1}{L{5.3cm}|}{\cellcolor[HTML]{FFFFFF} Identifikator des zu palettierenden Produkts und Schlüssel zur Auswahl produktspezifischer Parameter} 
		\\ \hline

		\multicolumn{1}{|L{2.5cm}|}{\cellcolor[HTML]{FFFFFF} BagQuantity}        	& 
		\multicolumn{1}{L{3.3cm}|}{\cellcolor[HTML]{FFFFFF} DIntServParam} 			& 
		\multicolumn{1}{L{3.17cm}|}{\cellcolor[HTML]{FFFFFF} 16\#2} 			& 
		\multicolumn{1}{L{5.3cm}|}{\cellcolor[HTML]{FFFFFF} Anzahl zu palettierender Säcke} 
		\\ \hline

		\multicolumn{1}{|L{2.5cm}|}{\cellcolor[HTML]{FFFFFF} PalletQuantity}        	& 
		\multicolumn{1}{L{3.3cm}|}{\cellcolor[HTML]{FFFFFF} DIntServParam} 			& 
		\multicolumn{1}{L{3.17cm}|}{\cellcolor[HTML]{FFFFFF} 16\#3} 			& 
		\multicolumn{1}{L{5.3cm}|}{\cellcolor[HTML]{FFFFFF} Anzahl zu palettierender Paletten} 
		\\ \hline

		\multicolumn{4}{|L{14.8cm}|}{\cellcolor[HTML]{E0E0E0}\textbf{Produkt- und verpackungsspezifische Parameter (Konfigurationsparameter)}} 
		\\ \hline

		\multicolumn{1}{|L{2.5cm}|}{\cellcolor[HTML]{E0E0E0}\textbf{Name}}        	& 
		\multicolumn{1}{L{3.3cm}|}{\cellcolor[HTML]{E0E0E0}\textbf{Schnittstellentyp}} 			& 
		\multicolumn{2}{L{8.65cm}|}{\cellcolor[HTML]{E0E0E0}\textbf{Beschreibung}} 
		\\ \hline

		\multicolumn{1}{|L{2.5cm}|}{\cellcolor[HTML]{FFFFFF} ProductDataSet}        	& 
		\multicolumn{1}{L{3.3cm}|}{\cellcolor[HTML]{FFFFFF} ArrayServParam of PalletizerProduct-DataStruct} 			& 
		\multicolumn{2}{L{8.65cm}|}{\cellcolor[HTML]{FFFFFF} Array produktspezifischer Parametersätze, enthält z.~B. das zu nutzende Packmuster} 
		\\ \hline

		\multicolumn{1}{|L{2.5cm}|}{\cellcolor[HTML]{FFFFFF} PackagingData-Set}        	& 
		\multicolumn{1}{L{3.3cm}|}{\cellcolor[HTML]{FFFFFF} ArrayServParam of PalletizerPackagingDataStruct} 			& 
		\multicolumn{2}{L{8.65cm}|}{\cellcolor[HTML]{FFFFFF} Array verpackungsspezifischer Parametersätze, enthält z.~B. die Abmessungen der zu verwendenden Paletten} 
		\\ \hline

		\multicolumn{4}{|L{14.8cm}|}{\cellcolor[HTML]{E0E0E0}\textbf{Konstruktionsspezifische Parameter (Konfigurationsparameter)}} 
		\\ \hline

		\multicolumn{4}{|L{14.8cm}|}{\cellcolor[HTML]{FFFFFF} Keine } 
		\\ \hline

		\multicolumn{4}{|L{14.8cm}|}{\cellcolor[HTML]{E0E0E0}\textbf{Reportwerte}} 
		\\ \hline

		\multicolumn{1}{|L{2.5cm}|}{\cellcolor[HTML]{E0E0E0}\textbf{Name}}        	& 
		\multicolumn{1}{L{3.3cm}|}{\cellcolor[HTML]{E0E0E0}\textbf{Schnittstellentyp}} 			& 
		\multicolumn{1}{L{3.17cm}|}{\cellcolor[HTML]{E0E0E0}\textbf{Prozedur}} 			& 
		\multicolumn{1}{L{5.3cm}|}{\cellcolor[HTML]{E0E0E0}\textbf{Beschreibung}} 
		\\ \hline

		\multicolumn{1}{|L{2.5cm}|}{\cellcolor[HTML]{FFFFFF} ProductName}        	& 
		\multicolumn{1}{L{3.3cm}|}{\cellcolor[HTML]{FFFFFF} StringView} 			& 
		\multicolumn{1}{L{3.17cm}|}{\cellcolor[HTML]{FFFFFF} 16\#1, 16\#2, 16\#3} 			& 
		\multicolumn{1}{L{5.3cm}|}{\cellcolor[HTML]{FFFFFF} Name des aktuell verpackten Produkttyps} 
		\\ \hline

		\multicolumn{1}{|L{2.5cm}|}{\cellcolor[HTML]{FFFFFF} PalletType-Name}        	& 
		\multicolumn{1}{L{3.3cm}|}{\cellcolor[HTML]{FFFFFF} StringView} 			& 
		\multicolumn{1}{L{3.17cm}|}{\cellcolor[HTML]{FFFFFF} 16\#1, 16\#2, 16\#3} 			& 
		\multicolumn{1}{L{5.3cm}|}{\cellcolor[HTML]{FFFFFF} Name des aktuell genutzten Palettentyps} 
		\\ \hline

		\multicolumn{1}{|L{2.5cm}|}{\cellcolor[HTML]{FFFFFF} BagCount}        	& 
		\multicolumn{1}{L{3.3cm}|}{\cellcolor[HTML]{FFFFFF} DIntView} 			& 
		\multicolumn{1}{L{3.17cm}|}{\cellcolor[HTML]{FFFFFF} 16\#1, 16\#2, 16\#3} 			& 
		\multicolumn{1}{L{5.3cm}|}{\cellcolor[HTML]{FFFFFF} Anzahl der bisher verpackten Säcke} 
		\\ \hline

		\multicolumn{1}{|L{2.5cm}|}{\cellcolor[HTML]{FFFFFF} PalletCount}        	& 
		\multicolumn{1}{L{3.3cm}|}{\cellcolor[HTML]{FFFFFF} DIntView} 			& 
		\multicolumn{1}{L{3.17cm}|}{\cellcolor[HTML]{FFFFFF} 16\#1, 16\#2, 16\#3} 			& 
		\multicolumn{1}{L{5.3cm}|}{\cellcolor[HTML]{FFFFFF} Anzahl der bisher verpackten Paletten} 
		\\ \hline

		\multicolumn{4}{|L{14.8cm}|}{\cellcolor[HTML]{E0E0E0}\textbf{Prozesswerte}} 
		\\ \hline

		\multicolumn{1}{|L{2.5cm}|}{\cellcolor[HTML]{E0E0E0}\textbf{Name}}        	& 
		\multicolumn{1}{L{3.3cm}|}{\cellcolor[HTML]{E0E0E0}\textbf{Schnittstellentyp}} 			& 
		\multicolumn{1}{L{3.17cm}|}{\cellcolor[HTML]{E0E0E0}\textbf{Prozedur}} 			& 
		\multicolumn{1}{L{5.3cm}|}{\cellcolor[HTML]{E0E0E0}\textbf{Beschreibung}} 
		\\ \hline

		\multicolumn{1}{|L{2.5cm}|}{\cellcolor[HTML]{FFFFFF} SuccessorClear-In}        	& 
		\multicolumn{1}{L{3.3cm}|}{\cellcolor[HTML]{FFFFFF} BinProcessValueIn} 			& 
		\multicolumn{1}{L{3.17cm}|}{\cellcolor[HTML]{FFFFFF} 16\#1, 16\#2, 16\#3} 			& 
		\multicolumn{1}{L{5.3cm}|}{\cellcolor[HTML]{FFFFFF} Eingang für das Freigabesignal von einer nachfolgenden LEA (relevant bei Nutzung der LEA in einer Logistics Line)} 
		\\ \hline

		\multicolumn{1}{|L{2.5cm}|}{\cellcolor[HTML]{FFFFFF} Predecessor-ClearOut}        	& 
		\multicolumn{1}{L{3.3cm}|}{\cellcolor[HTML]{FFFFFF} BinProcessValueIn} 			& 
		\multicolumn{1}{L{3.17cm}|}{\cellcolor[HTML]{FFFFFF} 16\#1, 16\#2, 16\#3} 			& 
		\multicolumn{1}{L{5.3cm}|}{\cellcolor[HTML]{FFFFFF} Ausgang für das Freigabesignal an eine vorhergehende LEA (relevant bei Nutzung der LEA in einer Logistics Line)} 
		\\ \hline

	 
	\end{longtable}
\end{footnotesize}

\subsubsection{Arbeitsweise}

Die Arbeitsweise des betrachteten PAL-Dienstes folgt der in Abschnitt~\ref{sec:CES} beschriebene Arbeitsweise eines CES und ist in Abbildung~\ref{fig:ArbeitsweisePal} dargestellt.

\begin{figure}[htbp]
	\centerline{\includegraphics[width=0.9\textwidth]{Inhalt/Abbildungen/99_Anhang/Spezifikation_LEA/Anwendung/Arbeitsweise_CES_PAL.png}}
	\caption{Arbeitsweise eines beispielhaften Palettierendienstes im CES-Betrieb}
	\label{fig:ArbeitsweisePal}
\end{figure}

Zu Beginn befindet sich der Dienst im Zustand IDLE. Es ist eine Vorbefüllung des \emph{ProductDataSet} und des \emph{PackagingDataSet} über die entsprechenden \emph{ArrayServParam}-Schnittstellen möglich. Bevor der Dienst gestartet wird, müssen zudem die gewünschte Prozedur gewählt und alle notwendigen auftragsspezifischen Parameter, d.~h. die \emph{ProductionId}, die \emph{ProductId}, die \emph{BagQuantity} (nur für Prozedur 16\#2) und die \emph{PalletQuantity} (nur für Prozedur 16\#3) eingestellt werden. Auf Basis der eingestellte \emph{ProductId} wird anschließend der zu nutzende produktspezifischen Parametersatz aus dem \emph{ProductDataSet} oder durch Abfrage beim LOL ermittelt und für die Bearbeitung ausgewählt (gemäß Abschnitt~\ref{subsec:Parametrierngsmech}). Mittels der im gewählten produktspezifischen Parametersatz enthaltenen \emph{PackagingId} wird auf gleiche Weise der zu nutzende verpackungsspezifische Parametersatz ermittelt und ausgewählt (gemäß Abschnitt~\ref{subsec:Parametrierngsmech}). Anschließend kann der PAL-Dienst gestartet werden und alle zum Auftrag gehörigen Säcke werden gemäß dem im Produktparametersatz hinterlegten Packmusters im Zustand EXECUTE palettiert. Der Palettiervorgang wird beendet, wenn die über \emph{BagQuantity} eingestellte Anzahl an Säcke erreicht ist (gilt für Prozedur 16\#2), wenn die über \emph{PalletQuantity} eingestellte Anzahl an Paletten erreicht ist (gilt für Prozedur 16\#3) oder wenn der Palettiervorgang durch einen \emph{Complete}-Befehl beendet wird (gilt für Prozedur 16\#1). Über einen \emph{Reset}-Befehl kann der Dienst in IDLE zurückgesetzt werden.

\subsubsection{Bedienbild}

Abbildung~\ref{fig:HmiPal} stellt das Bedienbild des PAL dar. Dieses ermöglicht die Eingabe der auftragsspezifischen Parameter, die Steuerung des LEA-Dienstes und die Anzeige der Reportwerte sowie Prozesswerte über dynamische Bedienbildobjekte gemäß \cite{PNO.2025.Part2}. Zudem enthält es ein Bild des PAL als statisches Bedienbildobjekt. Dieses Bild ist wie in Abbildung~\ref{fig:HmiPalMtpAnhang} dargestellt im Anhang des MTP des PAL hinterlegt. Es ist mit der ECLASS-Referenz \enquote{36521201} versehen, was der semantischen Referenz für einen Palettierer entspricht. Ein LOL, der diese LEA integriert, kann folglich entscheiden, ob er das Bild im MTP-Anhang oder ein potenziell in seiner Grafikbibliothek vorhandenes Bild eines Palettierers verwendet.

\begin{figure}[htbp]
	\centerline{\includegraphics[width=1.0\textwidth]{Inhalt/Abbildungen/99_Anhang/Spezifikation_LEA/Anwendung/HMI_PAL_PVs.png}}
	\caption{Bedienbild einer LEA zur Palettierung}
	\label{fig:HmiPal}
\end{figure}

\begin{figure}[htbp]
	\centerline{\includegraphics[width=0.5\textwidth]{Inhalt/Abbildungen/99_Anhang/Spezifikation_LEA/Anwendung/MTP_HMI_PAL.drawio.png}}
	\caption{Bilddatei einer LEA zur Palettierung innerhalb des MTP}
	\label{fig:HmiPalMtpAnhang}
\end{figure}

\subsection{Anwendungsbeispiel Stretchen-LEA}
\label{subsec:StretchenLea}

Zur Veranschaulichung der in Kapitel~\ref{chap:Art1LEA} beschriebenen Konzepte zeigt dieser Abschnitt die beispielhafte Anwendung dieser Konzepte auf ein Stretchen-LEA (SH), welches CES-basiert oder SES-basiert arbeiten kann.

\subsubsection{Dienstespezifikation}

Tabelle~\ref{tab:ShDienst} gibt einen Überblick über den Dienst des SH inkl. dessen Prozeduren, Parametern, Reportwerten und Prozesswerten.

% Steckbrief SH-Dienst
\begin{footnotesize}
	\begin{longtable}{|
		>{\columncolor[HTML]{E0E0E0}}l
		>{\columncolor[HTML]{E0E0E0}}l
		>{\columncolor[HTML]{E0E0E0}}l 
		>{\columncolor[HTML]{FFFFFF}}l |}
		\caption{Steckbrief eines beispielhaften Stretchendienstes}
		\label{tab:ShDienst}\\
		\hline

		\multicolumn{4}{|L{14.8cm}|}{\cellcolor[HTML]{E0E0E0}\textbf{Stretchen-Dienst}} 
		\\ \hline

		\multicolumn{4}{|L{14.8cm}|}{\cellcolor[HTML]{FFFFFF} Dieser Dienst kann abhängig von seiner Parametrierung Foliensicherungen auf Paletten mit Säcken oder Oktabins applizieren oder Folieninlays in leere Oktabins einlegen. Je nach Anlagenkontext kann das SH auftragsorientiert arbeiten oder bedarfsorientiert seine verschiedenen Aufgaben ggf. für verschiedene Verpackungsprozesse erfüllen.} 
		\\ \hline

		\multicolumn{4}{|L{14.8cm}|}{\cellcolor[HTML]{E0E0E0}\textbf{Prozeduren}} 
		\\ \hline

		\multicolumn{1}{|L{2.5cm}|}{\cellcolor[HTML]{E0E0E0}\textbf{Prozedur-ID}}        	& 
		\multicolumn{1}{L{3.3cm}|}{\cellcolor[HTML]{E0E0E0}\textbf{Name}} 			& 
		\multicolumn{2}{L{8.65cm}|}{\cellcolor[HTML]{E0E0E0}\textbf{Beschreibung}} 
		\\ \hline

		\multicolumn{1}{|L{2.5cm}|}{\cellcolor[HTML]{FFFFFF} 16\#1}                     & 
		\multicolumn{1}{L{3.3cm}|}{\cellcolor[HTML]{FFFFFF} CES\_Continuous (kontinuierlich)}              &                                              
		\multicolumn{2}{L{8.65cm}|}{\cellcolor[HTML]{FFFFFF} Kontinuierliches Ausführen \underline{einer} der o.~g. Aufgaben auf immer gleiche Weise gemäß Auftragsdaten}              
		\\ \hline

		\multicolumn{1}{|L{2.5cm}|}{\cellcolor[HTML]{FFFFFF} 16\#2}                     & 
		\multicolumn{1}{L{3.3cm}|}{\cellcolor[HTML]{FFFFFF} CES\_PalletQuantity (selbstbeendend)}              &                                              
		\multicolumn{2}{L{8.65cm}|}{\cellcolor[HTML]{FFFFFF} Ausführen \underline{einer} der o.~g. Aufgaben auf immer gleiche Weise an einer bestimmten Anzahl an LOs gemäß Auftragsdaten}              
		\\ \hline

		\multicolumn{1}{|L{2.5cm}|}{\cellcolor[HTML]{FFFFFF} 16\#3}                     & 
		\multicolumn{1}{L{3.3cm}|}{\cellcolor[HTML]{FFFFFF} SES\_Continuous (kontinuierlich)}              &                                              
		\multicolumn{2}{L{8.65cm}|}{\cellcolor[HTML]{FFFFFF} Flexible Ausführung der \underline{unterschiedlichen} o.~g. Aufgaben je nach ankommendem LO}              
		\\ \hline

		\multicolumn{4}{|L{14.8cm}|}{\cellcolor[HTML]{E0E0E0}\textbf{Auftragsspezifische Parameter (Prozedurparameter)}} 
		\\ \hline

		\multicolumn{1}{|L{2.5cm}|}{\cellcolor[HTML]{E0E0E0}\textbf{Name}}        	& 
		\multicolumn{1}{L{3.3cm}|}{\cellcolor[HTML]{E0E0E0}\textbf{Schnittstellentyp}} 			& 
		\multicolumn{1}{L{3.17cm}|}{\cellcolor[HTML]{E0E0E0}\textbf{Prozedur}} 			& 
		\multicolumn{1}{L{5.3cm}|}{\cellcolor[HTML]{E0E0E0}\textbf{Beschreibung}} 
		\\ \hline

		
		\multicolumn{1}{|L{2.5cm}|}{\cellcolor[HTML]{FFFFFF} ProductionId}        	& 
		\multicolumn{1}{L{3.3cm}|}{\cellcolor[HTML]{FFFFFF} StringServParam} 			& 
		\multicolumn{1}{L{3.17cm}|}{\cellcolor[HTML]{FFFFFF} 16\#1, 16\#2} 			& 
		\multicolumn{1}{L{5.3cm}|}{\cellcolor[HTML]{FFFFFF} Auftragsnummer} 
		\\ \hline

		\multicolumn{1}{|L{2.5cm}|}{\cellcolor[HTML]{FFFFFF} ProductId}        	& 
		\multicolumn{1}{L{3.3cm}|}{\cellcolor[HTML]{FFFFFF} DIntServParam} 			& 
		\multicolumn{1}{L{3.17cm}|}{\cellcolor[HTML]{FFFFFF} 16\#1, 16\#2} 			& 
		\multicolumn{1}{L{5.3cm}|}{\cellcolor[HTML]{FFFFFF} Identifikator des zu verarbeitenden Produkts und Schlüssel zur Auswahl produktspezifischer Parameter} 
		\\ \hline

		\multicolumn{1}{|L{2.5cm}|}{\cellcolor[HTML]{FFFFFF} PalletQuantity}        	& 
		\multicolumn{1}{L{3.3cm}|}{\cellcolor[HTML]{FFFFFF} DIntServParam} 			& 
		\multicolumn{1}{L{3.17cm}|}{\cellcolor[HTML]{FFFFFF} 16\#2} 			& 
		\multicolumn{1}{L{5.3cm}|}{\cellcolor[HTML]{FFFFFF} Anzahl zu verarbeitender LOs } 
		\\ \hline

		\multicolumn{4}{|L{14.8cm}|}{\cellcolor[HTML]{E0E0E0}\textbf{Produkt- und verpackungsspezifische Parameter (Konfigurationsparameter)}} 
		\\ \hline

		\multicolumn{1}{|L{2.5cm}|}{\cellcolor[HTML]{E0E0E0}\textbf{Name}}        	& 
		\multicolumn{1}{L{3.3cm}|}{\cellcolor[HTML]{E0E0E0}\textbf{Schnittstellentyp}} 			& 
		\multicolumn{2}{L{8.65cm}|}{\cellcolor[HTML]{E0E0E0}\textbf{Beschreibung}} 
		\\ \hline

		\multicolumn{1}{|L{2.5cm}|}{\cellcolor[HTML]{FFFFFF} ProductDataSet}        	& 
		\multicolumn{1}{L{3.3cm}|}{\cellcolor[HTML]{FFFFFF} ArrayServParam of StretchHoodProductDataStruct} 			& 
		\multicolumn{2}{L{8.65cm}|}{\cellcolor[HTML]{FFFFFF} Array produktspezifischer Parametersätze, enthält z.~B. die zu nutzenden Stretchparameter} 
		\\ \hline

		\multicolumn{1}{|L{2.5cm}|}{\cellcolor[HTML]{FFFFFF} PackagingData-Set}        	& 
		\multicolumn{1}{L{3.3cm}|}{\cellcolor[HTML]{FFFFFF} ArrayServParam of StretchHoodPackagingDataStruct} 			& 
		\multicolumn{2}{L{8.65cm}|}{\cellcolor[HTML]{FFFFFF} Array verpackungsspezifischer Parametersätze, enthält z.~B. die Merkmale der zu verwendenden Stretchfolie} 
		\\ \hline

		\multicolumn{4}{|L{14.8cm}|}{\cellcolor[HTML]{E0E0E0}\textbf{Konstruktionsspezifische Parameter (Konfigurationsparameter)}} 
		\\ \hline

		\multicolumn{1}{|L{2.5cm}|}{\cellcolor[HTML]{E0E0E0}\textbf{Name}}        	& 
		\multicolumn{1}{L{3.3cm}|}{\cellcolor[HTML]{E0E0E0}\textbf{Schnittstellentyp}} 			& 
		\multicolumn{2}{L{8.65cm}|}{\cellcolor[HTML]{E0E0E0}\textbf{Beschreibung}} 
		\\ \hline

		\multicolumn{1}{|L{2.5cm}|}{\cellcolor[HTML]{FFFFFF} FoilTypeId}        	& 
		\multicolumn{1}{L{3.3cm}|}{\cellcolor[HTML]{FFFFFF} DIntServParam} 			& 
		\multicolumn{2}{L{8.65cm}|}{\cellcolor[HTML]{FFFFFF} Identifikator des eingelegten Folientyps} 
		\\ \hline

		\multicolumn{4}{|L{14.8cm}|}{\cellcolor[HTML]{E0E0E0}\textbf{Reportwerte}} 
		\\ \hline

		\multicolumn{1}{|L{2.5cm}|}{\cellcolor[HTML]{E0E0E0}\textbf{Name}}        	& 
		\multicolumn{1}{L{3.3cm}|}{\cellcolor[HTML]{E0E0E0}\textbf{Schnittstellentyp}} 			& 
		\multicolumn{1}{L{3.17cm}|}{\cellcolor[HTML]{E0E0E0}\textbf{Prozedur}} 			& 
		\multicolumn{1}{L{5.3cm}|}{\cellcolor[HTML]{E0E0E0}\textbf{Beschreibung}} 
		\\ \hline

		\multicolumn{1}{|L{2.5cm}|}{\cellcolor[HTML]{FFFFFF} ProductName}        	& 
		\multicolumn{1}{L{3.3cm}|}{\cellcolor[HTML]{FFFFFF} StringView} 			& 
		\multicolumn{1}{L{3.17cm}|}{\cellcolor[HTML]{FFFFFF} 16\#1, 16\#2, 16\#3} 			& 
		\multicolumn{1}{L{5.3cm}|}{\cellcolor[HTML]{FFFFFF} Name des aktuell verpackten Produkttyps} 
		\\ \hline

		\multicolumn{1}{|L{2.5cm}|}{\cellcolor[HTML]{FFFFFF} FoilTypeName}        	& 
		\multicolumn{1}{L{3.3cm}|}{\cellcolor[HTML]{FFFFFF} StringView} 			& 
		\multicolumn{1}{L{3.17cm}|}{\cellcolor[HTML]{FFFFFF} 16\#1, 16\#2, 16\#3} 			& 
		\multicolumn{1}{L{5.3cm}|}{\cellcolor[HTML]{FFFFFF} Name des aktuell genutzten Folientyps} 
		\\ \hline

		\multicolumn{1}{|L{2.5cm}|}{\cellcolor[HTML]{FFFFFF} PalletCount}        	& 
		\multicolumn{1}{L{3.3cm}|}{\cellcolor[HTML]{FFFFFF} DIntView} 			& 
		\multicolumn{1}{L{3.17cm}|}{\cellcolor[HTML]{FFFFFF} 16\#1, 16\#2, 16\#3} 			& 
		\multicolumn{1}{L{5.3cm}|}{\cellcolor[HTML]{FFFFFF} Anzahl der bisher verarbeiteten LOs} 
		\\ \hline

		\multicolumn{4}{|L{14.8cm}|}{\cellcolor[HTML]{E0E0E0}\textbf{Prozesswerte}} 
		\\ \hline

		\multicolumn{1}{|L{2.5cm}|}{\cellcolor[HTML]{E0E0E0}\textbf{Name}}        	& 
		\multicolumn{1}{L{3.3cm}|}{\cellcolor[HTML]{E0E0E0}\textbf{Schnittstellentyp}} 			& 
		\multicolumn{1}{L{3.17cm}|}{\cellcolor[HTML]{E0E0E0}\textbf{Prozedur}} 			& 
		\multicolumn{1}{L{5.3cm}|}{\cellcolor[HTML]{E0E0E0}\textbf{Beschreibung}} 
		\\ \hline

		\multicolumn{1}{|L{2.5cm}|}{\cellcolor[HTML]{FFFFFF} SuccessorClear-In}        	& 
		\multicolumn{1}{L{3.3cm}|}{\cellcolor[HTML]{FFFFFF} BinProcessValueIn} 			& 
		\multicolumn{1}{L{3.17cm}|}{\cellcolor[HTML]{FFFFFF} 16\#1, 16\#2, 16\#3} 			& 
		\multicolumn{1}{L{5.3cm}|}{\cellcolor[HTML]{FFFFFF} Eingang für das Freigabesignal von einer nachfolgenden LEA (relevant bei Nutzung der LEA in einer Logistics Line)} 
		\\ \hline

		\multicolumn{1}{|L{2.5cm}|}{\cellcolor[HTML]{FFFFFF} Predecessor-ClearOut}        	& 
		\multicolumn{1}{L{3.3cm}|}{\cellcolor[HTML]{FFFFFF} BinProcessValueIn} 			& 
		\multicolumn{1}{L{3.17cm}|}{\cellcolor[HTML]{FFFFFF} 16\#1, 16\#2, 16\#3} 			& 
		\multicolumn{1}{L{5.3cm}|}{\cellcolor[HTML]{FFFFFF} Ausgang für das Freigabesignal an eine vorhergehende LEA (relevant bei Nutzung der LEA in einer Logistics Line)} 
		\\ \hline

	 
	\end{longtable}
\end{footnotesize}

\subsubsection{Arbeitsweise}

Die Arbeitsweise der beiden CES-Prozeduren ähnelt der Arbeitsweise des PAL-Dienstes. Daher wird in diesem Abschnitt die Arbeitsweise der SES-Prozedur des SH-Dienstes (16\#3) fokussiert. Diese folgt der in Abschnitt~\ref{sec:SES} beschriebene Arbeitsweise eines SES und ist in Abbildung~\ref{fig:ArbeitsweiseSh} dargestellt.

\begin{figure}[htbp]
	\centerline{\includegraphics[width=0.85\textwidth]{Inhalt/Abbildungen/99_Anhang/Spezifikation_LEA/Anwendung/Arbeitsweise_SES_SH.png}}
	\caption{Arbeitsweise eines beispielhaften Stretchendienstes im SES-Betrieb}
	\label{fig:ArbeitsweiseSh}
\end{figure}

Zu Beginn befindet sich der Dienst im Zustand IDLE. Es ist eine Vorbefüllung des \emph{ProductDataSet} und des \emph{PackagingDataSet} über die entsprechenden \emph{ArrayServParam}-Schnittstellen möglich. Zudem muss die LEA mit dem gewünschten Folientyp bestückt und dessen ID am konstruktionsspezifischen Parameter \emph{FoilTypeId} eingestellt werden. Mit dieser ID können die zugehörigen verpackungsspezifischen Parameter aus dem \emph{PackagingDataSet} ausgewählt oder vom Parameter-Management des LO angefragt werden (gemäß Abschnitt~\ref{subsec:Parametrierngsmech}).

Bevor der Dienst gestartet wird, muss zudem die gewünschte Prozedur gewählt werden (hier: 16\#3). Anschließend wird die Prozedur gestartet und wechselt in den Zustand PAUSED. Trifft ein LO ein, erfolgt ein Zustandswechsel nach RESUMING und die \emph{ProductId} sowie der \emph{LogisticsObjectStatus} des LO werden ermittelt (siehe auch Kapitel~\ref{chap:Art3LA}, \emph{ProductId} und \emph{LogisticsObjectStatus} als Teil der Schnittstelle eines Transportdienstes). Auf Basis dessen wird anschließend der zu nutzende produktspezifischen Parametersatz aus dem \emph{ProductDataSet} oder durch Abfrage beim LOL ermittelt und für die Bearbeitung ausgewählt (gemäß Abschnitt~\ref{subsec:Parametrierngsmech}). Es folgt ein Zustandswechsel nach EXECUTE. Als Grundlage für die Bearbeitung wird die im Produktdatensatz enthaltene \emph{PackagingId} mit der an der LEA eingestellten \emph{FoilTypeId} verglichen. Stimmen diese überein, wird das eingetroffene LO entsprechend der eingestellten Parameter bearbeitet. Stimmen diese nicht überein, kann die Palette vom LEA nicht bearbeitet werden und wird abgewiesen (siehe auch Abschnitt~\ref{subsec:Prozessmodell}, TransportDecline). 

Anschließend wechselt der Dienst zurück in den Zustand PAUSED und wartet auf das nächste ankommende LO. Sollen keine weiteren LOs mehr bearbeitet werden, kann der Dienst über einen \emph{Complete}-Befehl beendet und über einen \emph{Reset}-Befehl in IDLE zurückgesetzt werden.

\subsubsection{Bedienbild}

Abbildung~\ref{fig:HmiSh} stellt das Bedienbild des SH dar. Dieses ermöglicht die Eingabe der auftragsspezifischen Parameter, die Steuerung des LEA-Dienstes und die Anzeige der Reportwerte sowie Prozesswerte über dynamische Bedienbildobjekte gemäß \cite{PNO.2025.Part2}. Zudem enthält es ein Bild des SH als statisches Bedienbildobjekt. Dieses Bild ist wie in Abbildung~\ref{fig:HmiShMtpAnhang} dargestellt im Anhang des MTP des SH hinterlegt. Da keine passende ECLASS-Referenz für Stretchhaubenmaschinen existiert, wird die Referenz \enquote{90900001} verwendet. Ein LOL, der diese LEA integriert, muss folglich das im MTP-Anhang mitgelieferte Bild verwenden.

\begin{figure}[htbp]
	\centerline{\includegraphics[width=1.0\textwidth]{Inhalt/Abbildungen/99_Anhang/Spezifikation_LEA/Anwendung/HMI_SH_PVs.png}}
	\caption{Bedienbild einer LEA zum Stretchen und Folieneinlegen}
	\label{fig:HmiSh}
\end{figure}

\begin{figure}[htbp]
	\centerline{\includegraphics[width=0.5\textwidth]{Inhalt/Abbildungen/99_Anhang/Spezifikation_LEA/Anwendung/MTP_HMI_SH.drawio.png}}
	\caption{Bilddatei einer LEA zum Stretchen innerhalb des MTP}
	\label{fig:HmiShMtpAnhang}
\end{figure}


\subsection{Anwendungsbeispiel Sackabfüllinie}
\label{subsec:LLAnwendung}

Zur Veranschaulichung der in Kapitel~\ref{chap:Art2LL} beschriebenen Konzepte zeigt dieser Abschnitt die beispielhafte Anwendung des Choreografiekonzepts auf die Sackabfülllinie des unter Abschnitt~\ref{subsec:Anwendungsbeispiel} beschriebenen Anwendungsfalls.

\subsubsection{Einführung des Anwendungsfalls}

Die Sackabfüllinie besteht aus drei LEAs - einer FFS-Maschine (FFS), einem Transportband (CONV) und einem Palettierer (PAL), welche jeweils mit einem Dienst im CES-Betrieb gemäß Kapitel~\ref{chap:Art1LEA} automatisiert sind. Vereinfachend wird angenommen, dass diese Dienste ausschließlich die in Tabelle~\ref{tab:BspDiensteLogistiklinie} gezeigten Prozeduren, Parameter, Prozesswerte und Reportwerte aufweisen. 

% Dienste des Anwendungsbeispiels
\begin{footnotesize}
	\begin{longtable}{|
		>{\columncolor[HTML]{E0E0E0}}l 
		>{\columncolor[HTML]{E0E0E0}}l 
		>{\columncolor[HTML]{E0E0E0}}l 
		>{\columncolor[HTML]{E0E0E0}}l 
		>{\columncolor[HTML]{E0E0E0}}l |}
		\caption{Dienste der beispielhaften Logistics Line inkl. deren Prozeduren, Parameter und Prozesswerte}
		\label{tab:BspDiensteLogistiklinie}\\
		\hline 

		\multicolumn{1}{|L{2cm}|}{\cellcolor[HTML]{E0E0E0}\textbf{Dienst (LEA)}}        & 
		\multicolumn{1}{L{3.5cm}|}{\cellcolor[HTML]{E0E0E0}\textbf{Prozedur (ID)}}        & 
		\multicolumn{1}{L{2.3cm}|}{\cellcolor[HTML]{E0E0E0} \textbf{Parameter}} &
		\multicolumn{1}{L{3.5cm}|}{\cellcolor[HTML]{E0E0E0} \textbf{Prozesswert}} &
		\multicolumn{1}{L{2.3cm}|}{\cellcolor[HTML]{E0E0E0} \textbf{Reportwerte}}
		\\ \hline

        \multicolumn{1}{|L{2cm}|}{\cellcolor[HTML]{FFFFFF} BagFilling }        & 
		\multicolumn{1}{L{3.5cm}|}{\cellcolor[HTML]{FFFFFF} Continuous (16\#1)}        & 
		\multicolumn{1}{L{2.3cm}|}{\cellcolor[HTML]{FFFFFF}  	ProductId  } &
		\multicolumn{1}{L{3.5cm}|}{\cellcolor[HTML]{FFFFFF} SuccessorClearIn}  & 
		\multicolumn{1}{L{2.3cm}|}{\cellcolor[HTML]{FFFFFF} BagsCount}     
		\\ \cline{2-5}

		\multicolumn{1}{|L{2cm}|}{\cellcolor[HTML]{FFFFFF} (FFS)}     & 
		\multicolumn{1}{L{3.5cm}|}{\cellcolor[HTML]{FFFFFF} BagQuantity (16\#2)}        & 
		\multicolumn{1}{L{2.3cm}|}{\cellcolor[HTML]{FFFFFF}  	ProductId  } &
		\multicolumn{1}{L{3.5cm}|}{\cellcolor[HTML]{FFFFFF} SuccessorClearIn} & 
		\multicolumn{1}{L{2.3cm}|}{\cellcolor[HTML]{FFFFFF} BagsCount} 
		\\ \cline{3-3} 

		\multicolumn{1}{|L{2cm}|}{\cellcolor[HTML]{FFFFFF} }     & 
		\multicolumn{1}{L{3.5cm}|}{\cellcolor[HTML]{FFFFFF} }        & 
		\multicolumn{1}{L{2.3cm}|}{\cellcolor[HTML]{FFFFFF}  	BagQuantity  } &
		\multicolumn{1}{L{3.5cm}|}{\cellcolor[HTML]{FFFFFF} }    & 
		\multicolumn{1}{L{2.3cm}|}{\cellcolor[HTML]{FFFFFF} }     
		\\ \hline

		\multicolumn{1}{|L{2cm}|}{\cellcolor[HTML]{FFFFFF} Conveying }        & 
		\multicolumn{1}{L{3.5cm}|}{\cellcolor[HTML]{FFFFFF} Continuous (16\#1)}        & 
		\multicolumn{1}{L{2.3cm}|}{\cellcolor[HTML]{FFFFFF}  	Speed  } &
		\multicolumn{1}{L{3.5cm}|}{\cellcolor[HTML]{FFFFFF} PredecessorClearOut} & 
		\multicolumn{1}{L{2.3cm}|}{\cellcolor[HTML]{FFFFFF} - } 
		\\ \cline{4-4}

		\multicolumn{1}{|L{2cm}|}{\cellcolor[HTML]{FFFFFF} (CONV)}        & 
		\multicolumn{1}{L{3.5cm}|}{\cellcolor[HTML]{FFFFFF} }        & 
		\multicolumn{1}{L{2.3cm}|}{\cellcolor[HTML]{FFFFFF}  	 } &
		\multicolumn{1}{L{3.5cm}|}{\cellcolor[HTML]{FFFFFF} SuccessorClearIn} & 
		\multicolumn{1}{L{2.3cm}|}{\cellcolor[HTML]{FFFFFF} } 
		\\ \hline

		\multicolumn{1}{|L{2cm}|}{\cellcolor[HTML]{FFFFFF} Palletizing }        & 
		\multicolumn{1}{L{3.5cm}|}{\cellcolor[HTML]{FFFFFF} Continuous (16\#1)}        & 
		\multicolumn{1}{L{2.3cm}|}{\cellcolor[HTML]{FFFFFF}  	ProductId  } &
		\multicolumn{1}{L{3.5cm}|}{\cellcolor[HTML]{FFFFFF} PredecessorClearOut} & 
		\multicolumn{1}{L{2.3cm}|}{\cellcolor[HTML]{FFFFFF} BagsCount} 
		\\ \cline{5-5}

		\multicolumn{1}{|L{2cm}|}{\cellcolor[HTML]{FFFFFF} (PAL)}        & 
		\multicolumn{1}{L{3.5cm}|}{\cellcolor[HTML]{FFFFFF} }        & 
		\multicolumn{1}{L{2.3cm}|}{\cellcolor[HTML]{FFFFFF}  	  } &
		\multicolumn{1}{L{3.5cm}|}{\cellcolor[HTML]{FFFFFF} } & 
		\multicolumn{1}{L{2.3cm}|}{\cellcolor[HTML]{FFFFFF} PalletsCount} 
		\\ \cline{2-5}

		\multicolumn{1}{|L{2cm}|}{\cellcolor[HTML]{FFFFFF} }     & 
		\multicolumn{1}{L{3.5cm}|}{\cellcolor[HTML]{FFFFFF} BagQuantity (16\#2)}        & 
		\multicolumn{1}{L{2.3cm}|}{\cellcolor[HTML]{FFFFFF}  	ProductId  } &
		\multicolumn{1}{L{3.5cm}|}{\cellcolor[HTML]{FFFFFF} PredecessorClearOut} & 
		\multicolumn{1}{L{2.3cm}|}{\cellcolor[HTML]{FFFFFF} BagsCount} 
		\\ \cline{3-3} \cline{5-5}

		\multicolumn{1}{|L{2cm}|}{\cellcolor[HTML]{FFFFFF} }     & 
		\multicolumn{1}{L{3.5cm}|}{\cellcolor[HTML]{FFFFFF} }        & 
		\multicolumn{1}{L{2.3cm}|}{\cellcolor[HTML]{FFFFFF}  	BagQuantity  } &
		\multicolumn{1}{L{3.5cm}|}{\cellcolor[HTML]{FFFFFF} } & 
		\multicolumn{1}{L{2.3cm}|}{\cellcolor[HTML]{FFFFFF} PalletsCount} 
		\\ \cline{2-5}

		\multicolumn{1}{|L{2cm}|}{\cellcolor[HTML]{FFFFFF} }     & 
		\multicolumn{1}{L{3.5cm}|}{\cellcolor[HTML]{FFFFFF} PalletQuantity (16\#3)}        & 
		\multicolumn{1}{L{2.3cm}|}{\cellcolor[HTML]{FFFFFF}  	ProductId  } &
		\multicolumn{1}{L{3.5cm}|}{\cellcolor[HTML]{FFFFFF} PredecessorClearOut} & 
		\multicolumn{1}{L{2.3cm}|}{\cellcolor[HTML]{FFFFFF} BagsCount} 
		\\ \cline{3-3} \cline{5-5}

		\multicolumn{1}{|L{2cm}|}{\cellcolor[HTML]{FFFFFF} }     & 
		\multicolumn{1}{L{3.5cm}|}{\cellcolor[HTML]{FFFFFF} }        & 
		\multicolumn{1}{L{2.3cm}|}{\cellcolor[HTML]{FFFFFF}  	PalletQuantity  } &
		\multicolumn{1}{L{3.5cm}|}{\cellcolor[HTML]{FFFFFF} } & 
		\multicolumn{1}{L{2.3cm}|}{\cellcolor[HTML]{FFFFFF} PalletsCount} 
		\\ \hline
  
	\end{longtable}
\end{footnotesize}

Die Choreografiekonfiguration im Sinne der Konfigurierbaren Logik und Konfigurierbaren Kommunikation hängt wesentlich von dem gewünschten Verhalten der Logistics Line ab. Tabelle~\ref{tab:SzenarienLL} zeigt beispielhafte Anwendungsszenarien, die hier als Grundlage angenommen werden. In anderen Anwendungskontexten können andere Szenarien und somit eine andere Choreografiekonfiguration notwendig sein.

% Anwendungsszenarien der Sackabfülllinie
\begin{footnotesize}
	\begin{longtable}{|
		>{\columncolor[HTML]{E0E0E0}}l 
		>{\columncolor[HTML]{FFFFFF}}l |}
		\caption{Anwendungsszenarien der beispielhaften Sackabfülllinie}
		\label{tab:SzenarienLL}\\
		\hline

		\multicolumn{1}{|L{2cm}|}{\cellcolor[HTML]{E0E0E0}\textbf{Szenario}}        & 
		\multicolumn{1}{L{12.63cm}|}{\cellcolor[HTML]{E0E0E0} \textbf{Beschreibung}} 
		\\ \hline

        \multicolumn{1}{|L{2cm}|}{\cellcolor[HTML]{FFFFFF} Szenario 1 -\newline Anfahren}        & 
		\multicolumn{1}{L{12.63cm}|}{\cellcolor[HTML]{FFFFFF}  Die Linie muss in geordneter Weise von hinten nach vorne hochgefahren werden, d.~h. vom PAL zur FFS. Der Dienst einer LEA muss folglich gestartet werden, sobald der Dienst der nachfolgenden LEA läuft.} 
		\\ \hline

		\multicolumn{1}{|L{2cm}|}{\cellcolor[HTML]{FFFFFF} Szenario 2 -\newline Leerfahren}        & 
		\multicolumn{1}{L{12.63cm}|}{\cellcolor[HTML]{FFFFFF}  Die Linie muss in geordneter Weise von vorne nach hinten leergefahren werden, d.~h. von der FFS zum PAL. Der Dienst einer LEA muss folglich beendet werdem, sobald der Dienst der vorhergehenden LEA beendet ist. } 
		\\ \hline

		\multicolumn{1}{|L{2cm}|}{\cellcolor[HTML]{FFFFFF} Szenario 3 -\newline Exception Handling}        & 
		\multicolumn{1}{L{12.63cm}|}{\cellcolor[HTML]{FFFFFF}  Wenn in der Verpackungslinie ein Fehler auftritt, müssen alle LEAs in einen Fehlerzustand übergehen. So muss der Dienst einer LEA in einen Fehlerzustand wechseln, sobald ein anderer LEA-Dienst oder ein vorhandener LOL einen Fehlerzustand anzeigt.} 
		\\ \hline

		\multicolumn{1}{|L{2cm}|}{\cellcolor[HTML]{FFFFFF} Szenario 4 -\newline Unhold}        & 
		\multicolumn{1}{L{12.63cm}|}{\cellcolor[HTML]{FFFFFF}  Wenn der Fehler nicht kritisch ist (Dienstzustand: HELD), muss es möglich sein, die Logistics Line nach der Fehlerbehebung wieder anzufahren. Der Dienst einer LEA muss also wieder angefahren werden, sobald der Dienst der nachfolgenden LEA wieder läuft.} 
		\\ \hline

		\multicolumn{1}{|L{2cm}|}{\cellcolor[HTML]{FFFFFF} Szenario 5 -\newline Zurücksetzen}        & 
		\multicolumn{1}{L{12.63cm}|}{\cellcolor[HTML]{FFFFFF}  Wenn ein kritischer Fehler aufgetreten ist oder der Abfüllvorgang abgeschlossen wurde, muss es möglich sein, die Logistics Line in den Initialzustand IDLE zurückzusetzen. Der Dienst einer LEA muss also zurückgesetzt werden, sobald der Dienst der nachfolgenden LEA zurückgesetzt wurde.} 
		\\ \hline

       

	\end{longtable}
\end{footnotesize}

\subsubsection{Choreografiekonfiguration - Vertikale Integration}

\textbf{Lead Service:} Gemäß Abschnitt~\ref{subsec:LineSchnittstelle} stellt ein Lead Service eine Diensteschnittstelle für die Funktionalität einer Logistics Line bereit, der als Internal Lead oder External Lead ausgeführt sein kann. In diesem Fall wird ein External Lead gewählt, da keiner der drei beteiligten Dienste, die Line-Funktionalität sinnvoll abbilden kann. Bei Betrachtung beispielsweise des Szenarios 1 (Anfahren) fällt auf, dass der Impuls zum Anfahren beim PAL-Service gegeben werden muss, während die Rückmeldung zum erfolgreichen Anfahren beim FFS-Service vorliegt. Ein Internal Lead ist somit nicht geeignet. Der External Lead Service kann z.~B. neben dem \emph{Palletizing}-Dienst auf der Steuerung des PAL laufen, da diese heute meist die Leitsteuerungen von Logistiklinien darstellen, aber auch auf einer separaten Steuerung. Im Lead Service werden beispielhaft fünf generische Prozeduren vorgesehen, die entsprechend den Verhaltensregeln der Konfigurierbaren Logik zu verschiedenen Einstellungen der Prozeduren der Follower Services führen.

\textbf{Parameter:} Zudem sind die in Tabelle~\ref{tab:BspDiensteLogistiklinie} dargestellten Parameter einzustellen. Es ist sinnvoll, die Parameter an möglichst wenigen LEAs zu konfigurieren, um die Line-Schnittstelle zu vereinfachen. Im vorliegenden Fall werden alle Parameter des \emph{Palletizing}-Dienstes sowie der \emph{Speed}-Parameter des \emph{Conveying}-Dienstes an der Line-Schnittstelle vorgesehen. Die Parameter \emph{ProductId} und \emph{BagQuantity} werden choreografieintern vom PAL- an das FFS-LEA übergeben.

\textbf{Reportwerte:} Als Reportwerte der Logistics Line werden die Reportwerte \emph{BagsCount} und \emph{PalletsCount} des PAL vorgesehen, welche die Anzahl der fertiggestellten Säcke bzw. Paletten der Logisics Line bereitstellen.

\textbf{Prozesswerte:} Die vorhandenen Prozesswerte werden in diesem Fall ausschließlich choreografieintern verarbeitete und nicht an der Line-Schnittstelle bereitgestellt.

\textbf{Bedienbilder:} Das Line-Bedienbild setzt sich aus drei LEA-Bildern als CustomSymbols zusammen, die mittels des ContextReference-Mechanismus in Kombination mit dem CustomSymbols-Mechanismus aus den LEA-MTPs bezogen werden. Zudem werden für die \emph{ServiceControl}-Schnittstellen des Lead Service und der drei LEA-Dienste sowie für die Parameter und Reportwerte der Line-Schnittstelle dynamische \emph{VisualObjects} vorgesehen, die durch Referenzierung mittels \emph{RC HasExternalMtpContext} auf die zugehörigen \emph{DataAssemblies} in den angehängten LEA-MTPs verweisen. Abbildung~\ref{fig:LineBedienbild} stellt das Line-Bedienbild dar. Zusätzlich werden die drei Bedienbilder der einzelnen LEAs durch Referenzierung mittels \emph{ReferencedPicture} in die Bedienbildhierarchie der Logistics Line übernommen.
 
\begin{figure}[htbp]
	\centerline{\includegraphics[width=1.0\textwidth]{Inhalt/Abbildungen/07_Linen-Automatisierung/Line_STARTING_RV2.png}}
	\caption{Beispielhaftes Line Bedienbild einer Sackabfülllinie}
	\label{fig:LineBedienbild}
\end{figure}

Die zuvor beschriebenen Elemente der Line-Schnittstelle und des Line-Bedienbilds werden gemäß den in Abschnitt~\ref{sec:KonzeptComposedMTP} beschriebenen Mechanismen in einem Composed MTP abgebildet.

\subsubsection{Choreografiekonfiguration - Horizontale Integration}

Zwischen den drei LEAs der Logistics Line werden prozedurale, parametrierende und verriegelnde Relationen sowie Konstanten definiert.

\textbf{Prozedurale Relationen:} Abbildung~\ref{fig:ProzRelationen} zeigt einen beispielhaften Ausschnitt aus einer Matrix, die die prozeduralen Relationen zwischen den LEAs der Sackabfülllinie, im Speziellen die Relationen zwischen FFS und CONV, beschreibt.

\begin{figure}[htbp]
	\centerline{\includegraphics[width=0.7\textwidth]{Inhalt/Abbildungen/07_Linen-Automatisierung/Proz_Relationen.png}}
	\caption{Beispielhafte prozedurale Relationen einer Sackabfülllinie}
	\label{fig:ProzRelationen}
\end{figure}

Die Zeilen stellen die Systemvariablen dar, die der \emph{Conveying}-Dienst anderen Diensten zur Verfügung stellen kann. Die Spalten stellen die Choreografieschnittstelle des \emph{BagFilling}-Dienstes dar, die auf Basis der Variablen anderer LEA angesteuert werden kann. Die Markierungen innerhalb der Matrix repräsentieren die Verhaltensregeln der Konfigurierbaren Logik. So beschreibt die grün markierte Regel beispielsweise, dass der Zustand EXECUTE des \emph{Conveying}-Dienstes den Befehl \emph{Start} beim \emph{BagFilling}-Dienst auslöst. Diese Regel trägt zum geordneten Anfahren der Linie bei (Szenario 1). Bei einer solchen direkte Relation hat eine Systemvariable eines Dienstes eine unmittelbare Auswirkung auf die Schnittstelle eines anderen Dienstes. Beispielsweise im Falle des Szenario 3 \enquote{Fehlerfall} (Abbildung~\ref{fig:ProzRelationen}, rote Markierungen) kommt es darüber hinaus vor, dass zwei Systemvariablen auf dieselbe Variable der Choreografieschnittstelle einwirken. Hier ist eine Vorverarbeitung der Systemvariablen notwendig (Abbildung~\ref{fig:ProzRelationen}, \enquote{Vorverarbeitung}), deren Ergebnis an die Choreografieschnittstelle geschrieben wird. In diesem Fall ist diese Vorverarbeitung ein logisches ODER (\enquote{O}). Alle anderen arithmetischen und logischen Funktionen können ebenfalls für eine solche Vorverarbeitung verwendet werden. Wie hier am Beispiel des \emph{Conveying}- und \emph{BagFilling}-Dienstes gezeigt, lassen sich auch alle anderen Verhaltensregeln zwischen den drei Diensten der Sackabfülllinie und dem External Lead realisieren.

Zudem ist festzulegen, wie die Prozedureinstellung des External Leads an die LEA-Dienste weitergegeben werden soll. In diesem Beispiel wird beim Lead Service eine selbstbeendende Prozedur 1 zum Verpacken einer bestimmten Anzahl von Paletten eingestellt. Dies führt dazu, dass beim \emph{BagFilling}-Dienst die Prozedur \emph{BagQuantity}, beim \emph{Conveying}-Dienst die Prozedur \emph{Continuous} und beim \emph{Palletizing}-Dienst die Prozedur \emph{PalletQuantity} eingestellt wird.

\textbf{Parametrierende Relationen:} Abbildung~\ref{fig:ParamRelationen} zeigt einen beispielhaften Ausschnitt aus einer Matrix, die die parametrierenden Relationen zwischen den LEAs der Sackabfülllinie, im Speziellen die Relationen zwischen PAL und FFS, beschreibt.

\begin{figure}[htbp]
	\centerline{\includegraphics[width=0.4\textwidth]{Inhalt/Abbildungen/07_Linen-Automatisierung/Param_Relationen.png}}
	\caption{Beispielhafte parametrierende Relationen einer Sackabfülllinie}
	\label{fig:ParamRelationen}
\end{figure}

Wie zuvor im Abschnitt \enquote{Choreografiekonfiguration - Vertikale Integration} beschrieben, bezieht der \emph{BagFilling}-Dienst die von ihm benötigten Parameter (d.~h. die \emph{ProductId} und die \emph{BagQuantity}) aus dem PAL. Im vorliegenden Fall ist hierfür keine Vorverarbeitung notwendig.

\textbf{Verriegelnde Relationen:} Abbildung~\ref{fig:InterlockRelationen} zeigt einen beispielhaften Ausschnitt aus einer Matrix, die die verriegelnden Relationen zwischen den LEAs der Sackabfülllinie, im Speziellen die Relationen zwischen PAL und CONV, beschreibt.

\begin{figure}[htbp]
	\centerline{\includegraphics[width=0.5\textwidth]{Inhalt/Abbildungen/07_Linen-Automatisierung/Interlock_Relationen.png}}
	\caption{Beispielhafte verriegelnde Relationen einer Sackabfülllinie}
	\label{fig:InterlockRelationen}
\end{figure}

Im vorliegenden Anwendungsfall werden die vorhandenen Prozesswerte zur Übertragung der Freigabesignale als verriegelnde Relationen ohne Vorverarbeitung übertragen.

\textbf{Konstanten:} Abbildung~\ref{fig:ConstRelationen} zeigt einen beispielhaften Ausschnitt aus einer Matrix, die die Definition von Konstanten der Sackabfülllinie beschreibt.

\begin{figure}[htbp]
	\centerline{\includegraphics[width=0.85\textwidth]{Inhalt/Abbildungen/07_Linen-Automatisierung/Const_Relationen.png}}
	\caption{Beispielhafte Konstanten einer Sackabfülllinie}
	\label{fig:ConstRelationen}
\end{figure}

Die Möglichkeit, Konstanten zu definieren, wird im Fall der Sackabfülllinie verwendet, um die Betriebsart der LEA-Dienste gemäß Abschnitt~\ref{subsec:LineSchnittstelle} zu definieren. Bei den drei LEA-Dienste \emph{BagFilling}, \emph{Conveying} und \emph{Palletizing} wird die Variable \emph{Access Mode} auf einen konstanten Wert von \enquote{1} gesetzt. Dies interpretiert die LEA-Steuerung als die Betriebsart \emph{Automatic Internal}. Beim Lead Service wird der \emph{Access Mode} auf einen konstanten Wert von \enquote{2} gesetzt. Dies interpretiert die LEA-Steuerung als die Betriebsart \emph{Operator}. Darüber hinaus können Konstanten beispielsweise für das konstante Setzen von Parameter wie die \emph{Speed}-Variable des CONV verwendet werden. Weiterhin können Freigabesignale konstant auf TRUE gesetzt werden, wenn eine LEA, wie im Beispiel der PAL, am Ende der Logistics Line befindet und nicht auf die Freigabe einer nachfolgenden LEA warten muss.

\begin{graybox}
	\textbf{Best Practices zur Choreografiekonfiguration in Logistics Lines:}
	\begin{itemize}
		\item Das Entwurfsmuster Configurable Communication ermöglicht grundsätzlich den Austausch aller in einer LEA verfügbaren Informationen mit anderen LEAs einer choreografierten Logistics Line. Um eine herstellerneutrale Interaktion zwischen den LEAs zu ermöglichen, sollten allerdings, wo immer möglich, Variablen der standardisierten MTP-basierten LEA-Schnittstellen gemäß Kapitel~\ref{chap:Art1LEA} ausgetauscht werden.
		\item Eine Ausnahme bildet die Schnittstelle des External Lead. Hierfür werden gemäß \cite{Stutz.2026} \enquote{DoRequest- und DoneReply-Relationen} empfohlen, die die Zusammenführung von einem unterschiedlichen Anfang und Ende einer Choreografierelation in mehreren LEA-Diensten ermöglichen.
	\end{itemize}
\end{graybox}


\subsection{Anwendungsbeispiel Transport PAL-Label-SH}
\label{subsec:LAAnwendung}

Zur Veranschaulichung der in Kapitel~\ref{chap:Art3LA} beschriebenen Konzepte zeigt dieser Abschnitt die beispielhafte Anwendung des Transportkonzepts. Ähnlich wie in Kapitel~\ref{chap:Art3LA} wird als Beispiel ein Push-Transport vom PAL zur SH1 herangezogen, allerdings soll auf dem Weg zusätzlich ein Kennzeichner-LEA (LABEL) mit einem Processing-Knoten angefahren werden. Abbildung~\ref{fig:TransportBeispiel} stellt die LEAs und ihre Transportknoten dar. 

\begin{figure}[htbp]
	\centerline{\includegraphics[width=0.6\textwidth]{Inhalt/Abbildungen/08_Bereichs-Automatisierung/Transport_Anwendung/Beispiel.png}}
	\caption{Beispiel eines Push-Transports von einem Palettierer, über eine Kennzeichnungseinheit zu einer Stretchhaubenmaschine}
	\label{fig:TransportBeispiel}
\end{figure}

\subsubsection{0) LEA-Integration}

Als Grundlage für die Durchführung des Transportprozesses werden die MTPs der drei LEAs gemäß den Mechanismen aus Abschnitt~\ref{subsec:Transportanbindung} in das Transport-Management importiert, \emph{Transport-Proxies} generiert und an die LEA-Transportknoten angebunden. Abbildung~\ref{fig:IntegrationLea} zeigt den Prozess der Integration der drei LEAs, Abbildung~\ref{fig:IntegrationNode} zeigt den Prozess der Integration aller in den LEAs verfügbaren Transportknoten.\footnote{Die einzelnen TK-Proxies im Transport-Management sind hier aus Gründen der Übersichtlichkeit nicht dargestellt.}

\begin{figure}[htbp]
	\centerline{\includegraphics[width=0.8\textwidth]{Inhalt/Abbildungen/08_Bereichs-Automatisierung/Transport_Anwendung/00_Integration_LEA.png}}
	\caption{Integration von LEAs in ein Transport-Management}
	\label{fig:IntegrationLea}
\end{figure}

\begin{figure}[htbp]
	\centerline{\includegraphics[width=1.0\textwidth]{Inhalt/Abbildungen/08_Bereichs-Automatisierung/Transport_Anwendung/00_Integration_Node.png}}
	\caption{Integration aller Transportknoten von LEAs in ein Transport-Management}
	\label{fig:IntegrationNode}
\end{figure}

\subsubsection{1) Transport-Initiierung}

Abbildung~\ref{fig:TransportInitiierung} zeigt den Prozess der Transport-Initiierung.

\begin{figure}[htbp]
	\centerline{\includegraphics[width=1.0\textwidth]{Inhalt/Abbildungen/08_Bereichs-Automatisierung/Transport_Anwendung/01_Transport_Initiierung.png}}
	\caption{Initiierung eines Push-Transportauftrags}
	\label{fig:TransportInitiierung}
\end{figure}

Im Kontext des gewählten Beispiels besitzt ausschließlich der PAL einen Order-Knoten zur Initiierung von Transportaufträgen. Da diesem zu Beginn kein Transportdienst zugeordnet ist, wird im Transport-Management ein neuer Transportdienst (ID: transA) erzeugt und gemäß den Mechanismen aus Abschnitt~\ref{subsec:Transportanbindung} mit dem Order-Knoten (ID: 1) verbunden.

Im Beispiel weist der PAL einen Push-Transportbedarf auf und startet den Transportdienst in der Prozedur \emph{PushRequested}. Dabei wird der \emph{NextNode} auf den InOutbound-Knoten des PALs (ID: 2) gesetzt. Diese ID erhält der PAL durch Auslesen der \emph{ProxyId} in der \emph{TransportNodeManager}-Schnittstelle seines InOutbound-Knotens und stellt sie am Prozedurparameter \emph{NextNode} des Transportdienstes ein. 

Nach dem Start des Transportdienstes wird dessen OPC-UA-Verbindung zum Order-Knoten getrennt, er läuft allerdings im Transport-Management weiter. Für den nun freien Order-Knoten wird ein neuer Transportdienst (ID: transB) erzeugt und dem Order-Knoten zugeordnet.

\subsubsection{2) Fahrt zum Startknoten}

Abbildung~\ref{fig:FahrtStart} zeigt den Prozess der Fahrt zum Startknoten.

\begin{figure}[htbp]
	\centerline{\includegraphics[width=0.85\textwidth]{Inhalt/Abbildungen/08_Bereichs-Automatisierung/Transport_Anwendung/02_Transport_Fahrt_Start.png}}
	\caption{Fahrt zum Startknoten}
	\label{fig:FahrtStart}
\end{figure}

Der laufende Transportdienst (ID: transA) wird einem passenden FMS über dessen Adapter zur Verfügung gestellt. Dies kann je nach Schnittstelle des FMS auf Initiative des Transport-Managements oder des FMS geschehen. Sobald das FMS freie Kapazitäten hat, wählt es intern ein passendes AGV zur Abarbeitung des Transportprozesses aus und meldet die ID  (hier: agv3) an einen Adapter zurück. Dieser hinterlegt diese ID in der Variable \emph{ResourceId} des Transportdienstes. Sobald das FMS mit der Ausführung des Transportprozesses beginnt, startet es über seinen Adapter den Transportdienst in der Prozedur \emph{TransportRequested} neu. Anschließend fährt das ausgewählte AGV zum InOutbound-Knoten des PAL als Startknoten des Transportprozesses.

\subsubsection{3) Anbindung des Transportdienstes an den Startknoten}

Abbildung~\ref{fig:AnbindungTransportdienst} zeigt den Prozess der Anbindung des Transportdienstes an den Startknoten.

\begin{figure}[htbp]
	\centerline{\includegraphics[width=0.96\textwidth]{Inhalt/Abbildungen/08_Bereichs-Automatisierung/Transport_Anwendung/03_Anbindung_Transportdienst.png}}
	\caption{Anbindung des Transportdienstes an den Startknoten}
	\label{fig:AnbindungTransportdienst}
\end{figure}
 
Kurz bevor das AGV am InOutbound-Knoten des PAL ankommt, meldet das FMS seinem Adapter die bevorstehende Ankunft. Dieser initiiert in der Folge einen Prozedurwechsel nach \emph{TransportAwaitRequested} am Transportdienst. Ist der InOutbound-Knoten nicht durch einen anderen Transportdienst belegt, ordnet das Transport-Management dem \emph{Transport-Proxy} dieses Knotens den Transportdienst (ID: transA) gemäß Abschnitt~\ref{subsec:Transportanbindung} zu. Der PAL prüft anschließend anhand der Daten der Transportdienst-Schnittstelle (z.~B. \emph{ProductId}), ob er das transportierte Produkt verarbeiten kann. So kann er beispielsweise die \emph{ProductId} des Transportdienstes mit der \emph{ProductId} des Verpackungsauftrags vergleichen, den er aktuell bearbeitet. 

Je nach Prüfergebnis schaltet der PAL die Prozedur des Transportdienstes auf \emph{TransportAwaited} oder \emph{TransportDeclined} um. Bei einem negativen Prüfergebnis (\emph{TransportDeclined}) muss der nächste anzufahrende Transportknoten (Prozedurparameter \emph{NextNode}) neu definiert werden. Dazu kann die LEA mittels \emph{TransportNodeRequest} einen alternativen Knoten beim LOL oder beim Operator anfordern und den Transport entsprechend umleiten. Alternativ kann das FMS das AGV an eine sichere Position leiten und das Transport-Management versetzt den Transportdienst in den Zustand HELD, was die Notwendigkeit eines manuellen Eingreifens signalisiert.

Im vorliegenden Beispiel wird von einem positiven Prüfergebnis ausgegangen. In diesem Fall setzt der PAL die Prozedur des Transportdienstes auf \emph{TransportAwaited}, was der Freigabe für das das AGV entspricht, an den InOutbound-Knoten heranzufahren. Signalisiert das FMS, dass das AGV am Knoten angekommen ist, initiiert dessen Adapter einen Prozedurwechsel nach \emph{TransportArrived}.

\subsubsection{4) Beladung des AGV}

Abbildung~\ref{fig:BeladungAgv} zeigt den Prozess der Beladung des AGV.

\begin{figure}[htbp]
	\centerline{\includegraphics[width=0.95\textwidth]{Inhalt/Abbildungen/08_Bereichs-Automatisierung/Transport_Anwendung/04_Beladung_AGV.png}}
	\caption{Beladung des AGV}
	\label{fig:BeladungAgv}
\end{figure}

Sobald der InOutbound-Knoten des PAL den Status \emph{TransportArrived} erkennt, initiiert es eine Prozedurwechsel nach \emph{TransferFromLea} und die Übergabe des LO (hier: Palette) an das AGV beginnt. Sobald die Übernahme abgeschlossen ist, wird der nächste anzufahrende Transportknoten (hier: Processing-Knoten des LABEL-LEA) gemäß den Mechanismen aus Abschnitt~\ref{subsec:ErmittlungNextNode} ermittelt und am Prozedurparameter \emph{NextNode} des Transportdienstes eingestellt. Dann startet die LEA den Transportdienst in der Prozedur \emph{TransferFromLeaSucceeded} neu. Das erkennt das Transport-Management und entkoppelt den Transportdienst vom InOutbound-Knoten des PAL, woraufhin dieser wieder frei für andere Transportaufträge ist.

\subsubsection{5) Fahrt zum nächsten Knoten inkl. Umrouten}

Abbildung~\ref{fig:FahrtNextNode} zeigt den Prozess der Fahrt zum Folgeknoten inkl. der Möglichkeit des Umroutens.

\begin{figure}[htbp]
	\centerline{\includegraphics[width=0.85\textwidth]{Inhalt/Abbildungen/08_Bereichs-Automatisierung/Transport_Anwendung/05_Transport_Fahrt_NextNode_Rerouting.png}}
	\caption{Fahrt zum nächsten Knoten inkl. Umrouten}
	\label{fig:FahrtNextNode}
\end{figure}

Sobald der Adapter des FMS die Prozedur \emph{TransferFromLeaSucceeded} erkennt, gibt er den Auftrag zur Fahrt an den nächsten Transportknoten (hier: Processing-Knoten des LABEL-LEA) an das FMS weiter. Dieses meldet, wenn der Transport läuft, und der Adapter initiiert einen Prozedurwechsel nach \emph{Transport}.

Das Transport-Management überwacht die Verbindungen zu den \emph{TransportClientManager}-Schnittstellen der LEAs. Wird ein Verbindungsabbruch festgestellt, oder an der Variable \emph{LeaStateCur} der Schnittstelle wird ein Fehlerzustand festgestellt, kann ein Umrouten der Transportdienste initiiert werden, die zur betreffenden fehlerhaften LEA unterwegs sind. In diesem Fall wird LOL-intern initiiert durch das Transport-Management ein alternativer anzufahrender Transportknoten ermittelt und am Prozedurparameter \emph{NextNode} des Transportdienstes eingestellt. Anschließend startet das Transport-Management den Transportdienst in der Prozedur \emph{Rerouting} neu. Auf diese Weise erkennt der Adapter des FMS das Umrouten und gibt die Zieländerung an das FMS weiter. Wenn das FMS das neue Ziel übernommen hat, wird durch dessen Adapter ein Prozedurwechsel nach \emph{Transport} initiiert.

\subsubsection{6) Processing}

Abbildung~\ref{fig:Processing} zeigt den Prozess der Bearbeitung eines LO an einem Processing-Knoten.

\begin{figure}[htbp]
	\centerline{\includegraphics[width=1.0\textwidth]{Inhalt/Abbildungen/08_Bereichs-Automatisierung/Transport_Anwendung/06_Transport_Processing.png}}
	\caption{Processing eines Logistikobjekts}
	\label{fig:Processing}
\end{figure}
 
Handelt es sich beim angefahrenen Transportknoten um einen Processing-Knoten, wird der Transportdienst nach dem gleichen Prinzip wie im Schritt 3 beschrieben an diesen gekoppelt. Anschließend beginnt die Bearbeitung (Processing) des LO, was durch einen vom Transportknoten der LEA initiierten Prozedurwechsel nach \emph{Processing} signalisiert wird. Hierbei findet \underline{keine} Übergabe des LO auf die LEA statt, sondern es verbleibt auf dem AGV. Sobald die Bearbeitung abgeschlossen ist, wird der nächste anzufahrende Transportknoten bestimmt und am Prozedurparameter \emph{NextNode} des Transportdienstes eingestellt. Anschließend wird vom LEA ein Prozedurwechsel nach \emph{ProcessingSucceeded} durchgeführt. Das erkennt das Transport-Management und entkoppelt den Transportdienst vom Processing-Knoten. Anschließend können beliebig viele weitere Processing-Knoten auf die gleiche Weise angefahren werden. Die Fahrt zum nächsten Transportknoten findet nach dem gleichen Prinzip wie in Schritt 5 beschrieben statt.
 
\subsubsection{7) Entladung des AGV}

Abbildung~\ref{fig:EntladungAgv} zeigt den Prozess der Entladung des AGV.

\begin{figure}[htbp]
	\centerline{\includegraphics[width=1.0\textwidth]{Inhalt/Abbildungen/08_Bereichs-Automatisierung/Transport_Anwendung/07_Entladung_AGV.png}}
	\caption{Entladung des AGV}
	\label{fig:EntladungAgv}
\end{figure}

Wenn kein (weiterer) Processing-Knoten angefahren werden soll, wird der Transport zum letzten anzufahrenden Transportknoten (Zielknoten, hier: Inbound-Knoten der SH1) nach dem gleichen Prinzip wie in Schritt 5 koordiniert. Kurz vor Ankunft des AGV wird der Transportdienst nach dem gleichen Prinzip wie in in Schritt 3 beschrieben an den Transportknoten gekoppelt. Anschließend wird das LO vom AGV an das die SH1 übergeben, was durch einen vom Transportknoten der LEA initiierten Prozedurwechsel zu \emph{TransferToLea} signalisiert wird. Nach erfolgreicher Übergabe löst die LEA einen Prozedurwechsel zu \emph{TransferToLeaSucceeded} aus. Durch das Transport-Management wird der Transportdienst schließlich vom Transportknoten getrennt und mittels \emph{Complete}-Befehl beendet.

